\section{Introducción}
\label{sec:introduccion}

La confiabilidad de la literatura científica depende de la coherencia entre la pregunta de investigación, el diseño de recolección, el método de análisis y la interpretación de los resultados. La discusión contemporánea sobre reproducibilidad ha mostrado que una parte relevante de los problemas no surge solo de errores de cálculo, sino también de omisiones de reporte, uso rutinario de pruebas de hipótesis y conclusiones que exceden el alcance del diseño observado \citep{ioannidis2005,osc2015,schmidt2025}.

En particular, el \pvalue\ no resume magnitud de efecto, relevancia práctica ni precisión de la estimación. Por ello, una estrategia más adecuada para la toma de decisiones científicas consiste en priorizar estimación, intervalos de confianza (\IC), supuestos explícitos y métricas de precisión y exactitud, por ejemplo cobertura, sesgo y error cuadrático medio en estudios de simulación \citep{farcomeni2017,mansournia2021,darrigo2024}. En este artículo, \textit{falla inferencial crítica} refiere a prácticas u omisiones que invalidan la generalización a la población objetivo o inducen una interpretación incorrecta del efecto estimado.

Aunque el debate sobre reproducibilidad es global, su traducción a diagnósticos cuantitativos comparables es desigual entre regiones \citep{beigel2024}. En Sudamérica existe todavía una brecha para cuantificar, con métricas estandarizadas, la prevalencia de fallas críticas de diseño, recolección, análisis y reporte en publicaciones empíricas de acceso abierto. Esta ausencia limita la formulación de lineamientos específicos y evaluables para mejorar la calidad metodológica en áreas como salud, ciencias sociales e ingeniería. Este trabajo cubre esa brecha mediante un muestreo probabilístico multietápico de revistas y artículos, junto con la estimación ponderada de prevalencias e intervalos de confianza.

El objetivo de este estudio es doble: en primer lugar, recuperar y estimar la prevalencia estricta de fallas inferenciales en gran escala sobre artículos empíricos sudamericanos de acceso abierto, aplicando un sistema de extracción iterativa asistido por inteligencia artificial sobre cientos de cuerpos de texto; en segundo lugar, cuantificar, mediante simulación Monte Carlo, el impacto de incoherencias entre planificación muestral, mecanismo real de selección y validez inferencial, con especial atención a escenarios de muestreo por conveniencia.

Las contribuciones principales del estudio son las siguientes. En primer término, se definen operacionalmente las fallas inferenciales críticas ---en diferentes graduaciones de severidad--- que pueden auditarse de manera reproducible. En segundo término, se implementa y valida un canal automatizado de extracción JSON utilizando modelos de lenguaje masivos, permitiendo ampliar dramáticamente el marco de auditoría sin sacrificar profundidad de comprensión. En tercer lugar, se estima transversalmente la prevalencia empírica de fallas estrictas a nivel de artículo. En cuarto lugar, se evalúa por simulación la cobertura, el sesgo y el \RMSE\ bajo diseños probabilísticos frente a selección por conveniencia. Finalmente, se propone un conjunto de recomendaciones editoriales y de reporte orientadas a mejorar la coherencia entre diseño, análisis e interpretación.

El artículo se organiza de la siguiente manera. La Sección~\ref{sec:marcoteorico} resume los errores conceptuales que sustentan el protocolo de auditoría. La Sección~\ref{sec:metodologia} describe el diseño muestral, la estrategia de codificación y la simulación. La Sección~\ref{sec:resultados} presenta los resultados empíricos y los hallazgos del experimento Monte Carlo. Las implicancias metodológicas, las recomendaciones editoriales y las conclusiones se discuten de forma integrada en la Sección~\ref{sec:consideraciones}.
